\begin{problem}[第35届全国中学生物理竞赛复赛][X射线衍射与DNA结构]{94}
七、用波长为 $633\ nm$ 的激光水平照射竖直圆珠笔中的小弹簧，在距离弹簧 $4.2\ m$ 的光屏（与激光水平照射方向垂直）上形成衍射图像，如图a所示。其右图与1952年拍摄的首张DNA分子双螺旋结构X射线衍射图像（图b）十分相似。
\begin{subproblem}
利用图 a 右图中给出的尺寸信息，通过测量估算弹簧钢丝的直径 $d_{1}$ 、弹簧圈的半径 $R$ 和弹簧的螺距 $p$；
\end{subproblem}
\begin{subproblem}
图 b 是用波长为 $0.15\ nm$ 的平行 X 射线照射 DNA 分子样品后，在距离样品 $9.0\ cm$ 的照相底片上拍摄的。假设 DNA 分子与底片平行，且均与 X 射线照射方向垂直。根据图 b 中给出的尺寸信息，试估算 DNA 螺旋结构的半径 $R'$ 和螺距 $p'$ 。

说明：由光学原理可知，弹簧上两段互成角度的细铁丝的衍射、干涉图像与两条成同样角度、相同宽度的狭缝的衍射、干涉图像一致。
\end{subproblem}
\end{problem}

